\documentclass[11pt]{article}
\usepackage[margin=1in]{geometry}
\usepackage{amsmath, amssymb, amsthm}
\usepackage{mathtools}

\newtheorem{theorem}{Theorem}
\newtheorem{lemma}[theorem]{Lemma}
\newtheorem{corollary}[theorem]{Corollary}
\newtheorem{proposition}[theorem]{Proposition}
\theoremstyle{remark}
\newtheorem*{remark}{Remark}

\title{T4 (refuted): the slope-3 eigenvalue of $\Pi_q$ is not rational over $\mathbb{Q}(q)$}
\author{Clio}
\date{2026-04-29}

\begin{document}
\maketitle

\section{The conjecture and what is wrong with it}

For each shape $\lambda$ in the three ``atom-cube'' shapes
\[
\bigl\{\,(4,2)\text{ at }n=6,\ (4,3)\text{ at }n=7,\ (5,2)\text{ at }n=7\,\bigr\},
\]
the characteristic polynomial $g_\lambda(t) \in \mathbb{Z}[q,t]$ of $\Pi_q$ acting on
the irreducible Specht module $V_\lambda$ has been computed (Theorems T1, T2, T3
of 2026-04-27/28). The Newton polygon of $g_\lambda$ with respect to the
$\mathfrak{a}$-adic valuation, $\mathfrak{a} := q^2+4q+1$, has the following shape:
\begin{align*}
(4,2)/n=6 &: \text{slopes } -2, -1, 0,\ \text{each of length 1};\\
(4,3)/n=7 &: \text{slope } -3 \text{ (length 1)} \text{ and slope } 0 \text{ (length 3)};\\
(5,2)/n=7 &: \text{slope } -3 \text{ (length 1)} \text{ and slope } 0 \text{ (length 4)}.
\end{align*}

The (now-refuted) conjecture proposed:

\begin{quote}
\textbf{Conjecture T4.} For each shape $\lambda$, the unique eigenvalue
$r_\lambda$ of $\Pi_q | V_\lambda$ with positive $\mathfrak{a}$-adic valuation
lies in $\mathbb{Q}(q)$ and admits a closed form
\[
r_\lambda(q) = q^{a_\lambda}(1+q)^{b_\lambda}\mathfrak{a}(q)^{c_\lambda}\,P_\lambda(q),
\]
with $P_\lambda \in \mathbb{Z}[q]$ palindromic-positive.
\end{quote}

The justification offered was: \emph{the Newton polygon has a slope
$-c_\lambda$ segment of length 1; therefore $g_\lambda$ admits a linear factor
$g_1$ over $\mathbb{Q}(q)$, so its root is rational.}

\textbf{This justification is wrong.} The Newton-polygon factorisation
\[
g_\lambda(t) \;=\; g_1(t)\cdot g_0(t)
\]
takes place \emph{over the completion} $\mathbb{Q}(q)^{\mathfrak{a}}$, not
over $\mathbb{Q}(q)$ itself. The factor $g_1$ is linear over the completion,
so its root lies in $\mathbb{Q}(q)^{\mathfrak{a}}$, but \emph{a priori} the root may not lie
in the global field $\mathbb{Q}(q)$. The conjecture conflates ``length-1 segment''
with ``$\mathbb{Q}(q)$-linear factor''; only the former is automatic.

This note proves the conjecture is false in the strongest possible way:
$g_\lambda(t)$ is in fact \emph{irreducible} over $\mathbb{Q}(q)$ for every
one of the three shapes, so \emph{no} eigenvalue lies in $\mathbb{Q}(q)$.

\section{The theorem}

\begin{theorem}\label{thm:irr}
For each of the three shapes
$\lambda \in \{(4,2)/n=6,\ (4,3)/n=7,\ (5,2)/n=7\}$,
the characteristic polynomial $g_\lambda(t) \in \mathbb{Z}[q,t]$ of
$\Pi_q$ acting on the Specht module $V_\lambda$ is \emph{irreducible}
over $\mathbb{Q}(q)$.
\end{theorem}

\begin{corollary}\label{cor:no-rational}
None of the eigenvalues of $\Pi_q$ on $V_\lambda$ lies in $\mathbb{Q}(q)$.
In particular, the conjectured rational closed form $r_\lambda(q)$ does
not exist for any of the three atom-cube shapes.
\end{corollary}

The corollary is immediate: each eigenvalue of $\Pi_q | V_\lambda$ is a
root of $g_\lambda(t)$. Were any such root in $\mathbb{Q}(q)$, the
factor $(t - r) \in \mathbb{Q}(q)[t]$ would be a linear factor of
$g_\lambda$, contradicting Theorem~\ref{thm:irr}.

\section{Proof of Theorem~\ref{thm:irr}}

The key ingredient is the standard specialization criterion.

\begin{lemma}\label{lem:spec}
Let $g(q,t) \in \mathbb{Z}[q,t]$ be monic of positive degree in $t$.
If there exists $q_0 \in \mathbb{Z}$ such that the specialization
$g(q_0, t) \in \mathbb{Z}[t]$ is irreducible over $\mathbb{Q}$, then
$g(q,t)$ is irreducible over $\mathbb{Q}(q)$.
\end{lemma}

\begin{proof}
Suppose, for contradiction, that $g$ factors over $\mathbb{Q}(q)$ as
$g = f_1 \cdot f_2$ with $f_1, f_2 \in \mathbb{Q}(q)[t]$ both of positive
$t$-degree. Since $g$ is monic in $t$, by absorbing units we may take
$f_1, f_2$ each monic in $t$. By Gauss's lemma applied to
$\mathbb{Z}[q,t] = \mathbb{Z}[q][t]$, the monic factors $f_1, f_2$ in
fact lie in $\mathbb{Z}[q][t]$. Specializing at $q = q_0$, we obtain
\[
g(q_0, t) \;=\; f_1(q_0,t)\,f_2(q_0,t) \in \mathbb{Z}[t],
\]
where each factor is monic in $t$ of the same positive degree as $f_i$,
hence a non-trivial factorization over $\mathbb{Z}$ (and so over
$\mathbb{Q}$). This contradicts the assumed irreducibility of
$g(q_0, t)$ over $\mathbb{Q}$.
\end{proof}

\begin{proof}[Proof of Theorem~\ref{thm:irr}]
By Lemma~\ref{lem:spec}, it suffices to exhibit, for each shape
$\lambda$, a single value $q_0 \in \mathbb{Z}$ at which
$g_\lambda(q_0, t) \in \mathbb{Z}[t]$ is irreducible over $\mathbb{Q}$.

We take $q_0 = 2$ throughout. The polynomial $g_\lambda$ is determined
from the elementary symmetric functions $\sigma_1, \dots, \sigma_r$ of
the eigenvalues (where $r$ is the rank of the Hecke-invariant subspace,
the Specht-internal rank that selects the dimension of the image of
$\Pi_q$). These $\sigma_k(q)$ were computed exactly via Lagrange
interpolation in T1, T2, T3 (\emph{cf.}~the proof documents
\texttt{2026-04-27-newton-polygon-proof.tex} and
\texttt{2026-04-28-newton-polygon-n7-concentrated.tex}).

\medskip

\noindent\textbf{Shape $(4,2)$ at $n=6$ ($r=3$).} Specializing at $q_0 = 2$:
\[
g_{(4,2)}(2,t) = t^3 - 2{,}479{,}572\,t^2 + 726{,}098{,}626{,}656\,t - 11{,}766{,}474{,}745{,}339{,}392.
\]
Direct factorisation of this cubic over $\mathbb{Q}$ (or, equivalently,
over $\mathbb{Z}$ by Gauss) yields the polynomial as a single
irreducible factor. The rational root theorem rules out rational roots
since none of the divisors of $11{,}766{,}474{,}745{,}339{,}392$ is a
root, and a degree-3 polynomial over $\mathbb{Q}$ is reducible iff it
has a rational root. Verified computationally via SymPy's
\texttt{Poly($\cdot$, t, domain="QQ").factor\_list()}.

\medskip

\noindent\textbf{Shape $(4,3)$ at $n=7$ ($r=4$).} Specializing at $q_0=2$:
\begin{multline*}
g_{(4,3)}(2,t) = t^4 - 709{,}843{,}176\,t^3 + 63{,}282{,}432{,}370{,}649{,}472\,t^2 \\
- 469{,}939{,}687{,}270{,}795{,}269{,}181{,}440\,t \\
+ 191{,}646{,}477{,}786{,}342{,}535{,}098{,}958{,}872{,}576.
\end{multline*}
SymPy's \texttt{factor\_list} over $\mathbb{Q}$ returns a single
irreducible factor.

\medskip

\noindent\textbf{Shape $(5,2)$ at $n=7$ ($r=5$).} Specializing at $q_0=2$:
\begin{multline*}
g_{(5,2)}(2,t) = t^5 - 2{,}412{,}957{,}924\,t^4 + 998{,}432{,}827{,}755{,}810{,}336\,t^3 \\
- 70{,}028{,}605{,}693{,}778{,}159{,}853{,}609{,}984\,t^2 \\
+ 1{,}312{,}296{,}401{,}180{,}928{,}340{,}893{,}912{,}771{,}502{,}080\,t \\
- 2{,}004{,}689{,}847{,}756{,}036{,}486{,}677{,}534{,}364{,}715{,}630{,}460{,}928.
\end{multline*}
SymPy's \texttt{factor\_list} over $\mathbb{Q}$ returns a single
irreducible factor.

\medskip

In each case the specialization $g_\lambda(2, t)$ is irreducible over
$\mathbb{Q}$, so by Lemma~\ref{lem:spec} the polynomial
$g_\lambda(q, t)$ is irreducible over $\mathbb{Q}(q)$.
\end{proof}

\begin{remark}[Cross-check]
We further verified irreducibility of $g_\lambda(q_0, t)$ at
$q_0 = 3$ for all three shapes and at $q_0 = 5$ for $(4,2)$.
Both factorisations again returned a single irreducible factor, in line
with the fact that the irreducibility statement at one specialization
already implies irreducibility globally. Independent SymPy
factorisations over $\mathbb{Q}(q)[t]$ directly (using
\texttt{Poly($\cdot$, t, domain=$\mathbb{Q}(q)$).factor\_list()})
confirm the same: each $g_\lambda$ remains a single factor of full
degree.
\end{remark}

\begin{remark}[Galois group]
Theorem~\ref{thm:irr} shows the Galois group $\mathrm{Gal}(L_\lambda /
\mathbb{Q}(q))$ of the splitting field $L_\lambda$ acts transitively
on the $r_\lambda$ eigenvalues. By the Newton-polygon information, the
decomposition group at the place $\mathfrak{a}$ has orbits of sizes
$(1,1,1)$, $(1,3)$, $(1,4)$ on the eigenvalues for the three shapes
respectively, since the Newton-polygon factorisation over the
completion $\mathbb{Q}(q)^{\mathfrak{a}}$ has factors of those
$t$-degrees. Determining the full Galois groups is not pursued here;
the relevant computational handle is the discriminant of $g_\lambda$
(easily extracted from the explicit polynomial).
\end{remark}

\section{Where the original argument breaks down}

The error in the original justification of T4 is the conflation of
two distinct factorisations:

\begin{itemize}
\item \textbf{Newton-polygon factorisation over $\mathbb{Q}(q)^\mathfrak{a}$.}
  Over the $\mathfrak{a}$-adic completion of $\mathbb{Q}(q)$, the
  polynomial $g_\lambda$ splits as
  \[
  g_\lambda(t) \;=\; g_1(t) \cdot g_0(t),
  \]
  where $g_1$ has all roots of $\mathfrak{a}$-adic valuation
  $c_\lambda = 3$ (resp.\ $2$, $1$, $0$ for the three slopes of $(4,2)$),
  $g_0$ has all roots of valuation $0$, and $\deg g_1$ equals the
  length of the corresponding slope. This is automatic from the
  polygon and is used in the proofs of T1, T2, T3.

\item \textbf{Polynomial factorisation over $\mathbb{Q}(q)$.}
  This is the question of whether $g_\lambda \in \mathbb{Q}(q)[t]$
  factors. Theorem~\ref{thm:irr} says it does not.
\end{itemize}

These two factorisations agree exactly when $\mathbb{Q}(q)$ is
\emph{Henselian} at the place $\mathfrak{a}$, equivalently when
$\mathbb{Q}(q)$ equals its $\mathfrak{a}$-adic completion --- which
is far from true. The atom $\mathfrak{a} = q^2+4q+1$ is irreducible of
degree $2$ in $\mathbb{Q}[q]$ (its discriminant is $12$, not a square),
so the residue field $\mathbb{Q}[q]/(\mathfrak{a}) \cong \mathbb{Q}(\sqrt{3})$
is a non-trivial extension of $\mathbb{Q}$. The completion
$\mathbb{Q}(q)^\mathfrak{a}$ contains, but strictly,
$\mathbb{Q}(q)$.

\section{What survives}

\begin{proposition}[Atom-adic existence and uniqueness, verified]
For each shape $\lambda$, there is a \emph{unique} root $r_\lambda$ of
$g_\lambda$ in the completion $\mathbb{Q}(q)^\mathfrak{a}$ with positive
$\mathfrak{a}$-adic valuation, and that valuation equals $c_\lambda$
(the highest atom-exponent dictated by the polygon: $2$, $3$, $3$ for
$(4,2)$, $(4,3)$, $(5,2)$ respectively). This root lies in
$\mathbb{Q}(q)^\mathfrak{a}$ but \emph{not} in $\mathbb{Q}(q)$.
\end{proposition}

This was the genuine content extractable from the Newton polygon
argument; only the \emph{rationality} claim was unwarranted.

\begin{proposition}[Leading coefficient in residue field]
Write $r_\lambda = \mathfrak{a}^{c_\lambda} \cdot u_\lambda$ with
$u_\lambda$ a unit in the local ring of $\mathbb{Q}(q)^\mathfrak{a}$.
Reducing modulo $\mathfrak{a}$ gives a leading coefficient
$\overline{u}_\lambda \in \mathbb{Q}(\alpha)$, where $\alpha$ is a
root of $\mathfrak{a}$, i.e.\ $\alpha = -2 + \sqrt{3}$.

A direct computation (from the leading Hensel approximation
$r_\lambda \equiv \sigma_r/\sigma_{r-1} \pmod{\mathfrak{a}^{c_\lambda + 1}}$,
which is valid since $\mathrm{ord}_\mathfrak{a}(g(\sigma_r/\sigma_{r-1}))
= c_\lambda + 1$ in each case) yields:
\begin{align*}
\overline{u}_{(4,2)} &= \tfrac{-850720 - 290912\,\alpha}{68361}
  \in \mathbb{Q}(\alpha),\\
\overline{u}_{(4,3)} &= \tfrac{78240950912 + 28791367552\,\alpha}{7949071559}
  \in \mathbb{Q}(\alpha),\\
\overline{u}_{(5,2)} &= \tfrac{-2281470720 + 1742825728\,\alpha}{1128105863}
  \in \mathbb{Q}(\alpha).
\end{align*}
The denominators arise from rationalising
$\overline{P_{r-1}(\alpha)}^{-1}$, where $P_{r-1}$ is the
remainder polynomial appearing in $\sigma_{r-1} = q^a (1+q)^b \cdot
P_{r-1}(q)$ with $P_{r-1}$ coprime to $\mathfrak{a}$.

These $\overline{u}_\lambda$ are not in $\mathbb{Q}$ (they have nonzero
$\alpha$-coefficient), confirming once more that $r_\lambda$ is not
$\mathbb{Q}(q)$-rational.
\end{proposition}

The cross-shape pattern from $A_9 = q^6 + 12q^5 + 52q^4 + 88q^3 + 52q^2 +
12q + 1$ (which appears as the palindromic core of $\sigma_r$ for
\emph{both} $(4,3)$ and $(5,2)$ at $n=7$) does \emph{not} translate
into a clean cross-shape relation between $\overline{u}_{(4,3)}$ and
$\overline{u}_{(5,2)}$: the denominators $7{,}949{,}071{,}559$ and
$1{,}128{,}105{,}863$ have no obvious common structure.

\section{Computational verification}

\textbf{Code.} \texttt{\~/projects/scratch/2026-04-29-rational-atom-eigenvalue/}
\begin{itemize}
\item \texttt{compute\_eigenvalue.py} — interpolates $\sigma_k(q)$,
  forms $g_\lambda(t)$, factors over $\mathbb{Q}(q)[t]$ via SymPy.
\item \texttt{specialize\_check.py} — specializes
  $g_\lambda(q_0, t)$ at $q_0 \in \{2, 3, 5\}$, factors over $\mathbb{Q}[t]$.
\item \texttt{check\_extensions.py} — factors over
  $\mathbb{Q}(\sqrt{3})(q)[t]$ (residue-field extension).
\item \texttt{leading\_unit.py} — extracts the leading residue-field
  coefficient $\overline{u}_\lambda \in \mathbb{Q}(\alpha)$.
\item \texttt{verify\_irreducibility.py} — Newton iteration showing
  $\sigma_r/\sigma_{r-1}$ is not a root.
\end{itemize}

\textbf{Verifications run.}
\begin{enumerate}
\item Each $\sigma_k(q) \in \mathbb{Z}[q]$ Lagrange-interpolated from
  $89$/$50$/$110$ integer evaluations (overdetermined relative to
  the degree bound; cross-checked).
\item $g_\lambda(t)$ over $\mathbb{Q}(q)[t]$:
  \texttt{Poly($\cdot$, t, domain=QQ(q)).factor\_list()} returns a
  single factor of full $t$-degree for each shape.
\item $g_\lambda(q_0, t)$ over $\mathbb{Q}[t]$ for
  $q_0 = 2, 3, 5$: irreducible in every case.
\item $g_\lambda(t)$ over $\mathbb{Q}(\sqrt{3})(q)[t]$: still
  irreducible. So the slope-$c_\lambda$ root does not lie even in the
  unramified extension corresponding to the residue field
  $\mathbb{Q}(\alpha)$.
\end{enumerate}

\section{Status}

T4 (the rational atom-eigenvalue conjecture) is \textbf{refuted}.

\begin{itemize}
\item The eigenvalues of $\Pi_q$ on the three atom-cube Specht modules
  are not individually in $\mathbb{Q}(q)$.
\item They generate a degree-$r$ extension of $\mathbb{Q}(q)$ (namely
  $\mathbb{Q}(q)[t]/(g_\lambda)$).
\item The original Newton-polygon argument was misapplied: the polygon
  factorisation lives over the completion, not the global field.
\item What is preserved: the symmetric functions $\sigma_k$ (the
  Eigenvalue Rosetta Stone factorisation $\sigma_k = q^a(1+q)^b
  \mathfrak{a}^c P_k(q)$ with $P_k$ palindromic) — these are
  $\mathbb{Q}(q)$-rational by construction. The atom-exponent grading
  of the eigenvalue \emph{multiset} (T1, T2, T3) is unchanged.
\end{itemize}

The $\sigma_k$-G1 program (palindromic-positivity at level
$\sigma_k$) is unaffected: it is a statement about the symmetric
functions, not about individual eigenvalues. So T4 was a \emph{stretch
goal} in the original PROVE.md (``what is the closed form of the
atom-cube eigenvalue?''), and the answer is now: there is no
$\mathbb{Q}(q)$-closed form, because no such closed form exists.

\medskip

\noindent\textbf{Open directions.}
\begin{itemize}
\item \textbf{Galois group.} Compute $\mathrm{disc}(g_\lambda)$ and
  determine $\mathrm{Gal}(g_\lambda/\mathbb{Q}(q))$. For $(4,2)$ this
  is a transitive subgroup of $S_3$ — either $S_3$ or $A_3 = \mathbb{Z}/3$.
  For $(4,3)$ it is a transitive subgroup of $S_4$. For $(5,2)$ it is
  a transitive subgroup of $S_5$.
\item \textbf{Higher-order Hensel coefficients.} The expansion
  $r_\lambda = \mathfrak{a}^{c_\lambda} \overline{u}_\lambda + \mathfrak{a}^{c_\lambda+1} \overline{u}^{(1)}_\lambda + \cdots$
  has each coefficient in $\mathbb{Q}(\alpha)$. Whether these
  coefficients have any pattern (e.g.\ all in $\mathbb{Q}$ from some
  point on, or following a recursion) is a natural follow-up.
\item \textbf{Why $g$ is irreducible.} What categorical / Hecke
  structural reason forces the eigenvalues of $\Pi_q$ on $V_\lambda$ to
  be Galois-conjugate over $\mathbb{Q}(q)$ for these three shapes? This
  is the genuinely interesting question that emerges from the
  refutation.
\end{itemize}

\end{document}
